% Section 6.5, from lectures/L06/README.md: what you should be able to do, questions to test
% yourself, and what comes after the course.

\needspace{10\baselineskip}
\section{Review}
\label{c6:sec:review}

\subsection*{What you should be able to do now}
\begin{itemize}
  \item Say what a watchdog detects, and name three failures it does not.
  \item Write the timed sequence correctly, and say what the hardware does when you do not.
  \item Choose a timeout from a worst-case loop time, and say why the choice is not obvious in
    either direction.
  \item Explain why \code{WDRF} must be cleared before \code{WDE} can be, and what goes wrong on the
    second run of a program that forgets.
  \item Choose a sleep mode from what the program needs left running, and name the two wake sources
    that survive every mode.
  \item State the AVR C calling convention: argument registers, return registers, call-saved
    registers, and the register the compiler keeps at zero.
  \item Call an assembly subroutine from C, and a C function from assembly, and say what each side
    is promising the other.
  \item Read \code{avr-gcc -S} output and compare it against assembly you wrote for the same job.
\end{itemize}

\needspace{6\baselineskip}
\subsection*{Questions to test yourself}
\begin{enumerate}
  \item Your program deadlocks waiting for an interrupt that never comes. Does the watchdog help?
    What if instead it computes a wrong answer very quickly?
  \item The two writes of the timed sequence must be within four cycles. Why is \code{cli} before
    them not optional?
  \item You write the byte \code{0x08} into \code{WDTCSR} meaning ``the four second timeout''. What
    did you actually ask for, and what happens next?
  \item A program is reset by the watchdog, restarts, and tries to switch the watchdog off. Why
    might it fail, and what has to happen first?
  \item You put the device into power-down and expect Timer1 to wake it in a second. What actually
    happens?
  \item You pass a 16-bit value to an assembly subroutine from C. Which registers does it arrive in,
    and which half is in the lower-numbered one?
\end{enumerate}

\needspace{6\baselineskip}
\subsection*{After the book}
This is the last chapter. What you have built is a driver library in assembly, and the habit that
goes with it: \textbf{compute the number first, then measure it, and take the disagreement
seriously.}

The companion repository has two written papers in \file{exam/}, if you want to test yourself on
paper. They gate nothing and nothing in this book requires them.

\textbf{Real-Time Kernel Design} is the course that follows this one: the same part, the same
harness, the same check by hand and by measurement, building a preemptive kernel one context switch
at a time. It names this material a prerequisite and it means it. Chapter~3's ISR contract,
Chapter~4's stack arithmetic and its loadable stack pointer (\secref{c4:sec:loadsp}), Chapter~5's
timer arithmetic, and this chapter's C ABI in both directions are used there rather than recapped.
The half of a kernel that must be assembly is about two files out of a dozen, and knowing
\emph{which} two is most of what that course is for.
